\documentclass[12pt]{article}

\usepackage{fontawesome}
\usepackage{hyperref}
\usepackage{xurl}
\usepackage{float}
\usepackage[tagged, highstructure]{accessibility}
\usepackage{bookmark}
\usepackage{enumitem}

\hypersetup{
    colorlinks=false,
    pdfborder={0 0 0},
    pdftitle={Final Presentation},
    pdfauthor={Adrianna Holden-Gouveia},
    pdfsubject={Information Security - Lab Assignment},
    pdflang={en-US},
    pdfkeywords={final presentation, security training, video, CRAAP, Brightspace}
}

% Configure PDF bookmarks for navigation
\bookmarksetup{
    numbered,
    open,
}

% Configure list spacing for better accessibility
\setlist{nosep}



\title{Final Presentation}
\author{
        Adrianna Holden-Gouveia \\
        Website: \href{https://aholdengouveia.name}{aholdengouveia.name}\\ 
        \date{\vspace{-5ex}}
        %Email: \href{mailto:admin@aholdengouveia.name}{admin@aholdengouveia.name} \\
%        \faLinkedin{: aholdengouveia} \\
        \faGithub {: aholdengouveia} \\
        \faTwitter {: aholdengouveia} \\
        }

%S\date{\today}


\begin{document}

\maketitle

\section*{Accessibility Notice}
This document is also available in HTML format at:

Link: \href{https://aholdengouveia.name/InfoSec/labs/finalpresentation.html}{\textbf{\nolinkurl{https://aholdengouveia.name/InfoSec/labs/finalpresentation.html}}}

The HTML version provides enhanced accessibility features including keyboard navigation, screen reader support, responsive design, dark mode support, and high contrast options.

\section*{Learning Outcomes}
By the end of this final presentation, you will be able to:
\begin{itemize}
    \item Research an information security topic that wasn't covered in class, using reliable sources checked with the CRAAP test
    \item Plan a security lesson for a specific audience, using language and examples that fit them
    \item Present a recorded lesson with a clear introduction, supporting points, conclusion, and visuals, within a time limit
    \item Give specific, constructive feedback on other students' presentations
\end{itemize}

\section*{Final Presentation Instructions}
Pretend you've been asked to give a short security training. Pick a topic from my Final Presentation Topics list, linked under Resources, pick who you're training, and prepare a 9 to 11 minute lesson. You'll record it, and share it with the class on Brightspace.

The topic has to be new. Don't repeat a topic we covered in a lab, such as malware, social engineering, firewalls, passwords, web security, physical security, data breaches, cloud security, forensics, or cleaning up infections. Your topic and presentation must be about information security.

\subsection*{Pick your audience}
Choose who you're training, and say who it is at the start of your presentation. For example: new employees at a small business, high school students, older family members, a sports team, or other college students. Your audience changes your examples, the words you use, and the steps you give them.

\subsection*{What your lesson must include}
\begin{itemize}
    \item Why this topic matters to your audience
    \item What the threat or problem is, and how it works, explained for your audience
    \item At least one real example, such as a news story or a case from the last few years
    \item What your audience can do: steps, tools, or habits that protect them
    \item Where they can go to learn more
    \item A clear introduction and conclusion
\end{itemize}

\subsection*{Recording rules}
\begin{itemize}
    \item Your video must be 9 to 11 minutes long.
    \item \textbf{It must be you presenting, in your own voice.} You don't have to show your face. AI voices, AI-generated presenters, and AI-written scripts aren't allowed.
    \item \textbf{Record with a screen capture}, so your slides and anything you show are on screen while you talk. Good free options: OBS Studio, a Zoom meeting where you share your screen and record, or the Screenshot app on a Mac (press Shift, Command, and 5).
    \item \textbf{Don't use PowerPoint's built-in recording}, such as Record Slide Show, or exporting your slides as a narrated video. These won't be accepted.
    \item Don't play other people's videos as part of your presentation. The presentation has to be you.
\end{itemize}

\section*{Steps}
Due dates and the grading rubric are in Brightspace.
\begin{enumerate}
    \item \textbf{Claim your topic.} In the Final Presentation discussion on Brightspace, start a thread. Make the title your name and your topic. In the body, write who your audience is, a short summary of what you'll cover, and your 5 sources with full links. A topic is only claimed when your thread has all of these. Threads can't be edited, and topics are first come, first served, going by Brightspace's time stamp. If someone already has your topic, pick another one.
    \item \textbf{Record and post your video.} Reply to your own thread with your video, so everyone can watch it and comment.
    \item \textbf{Comment on everyone else's video.} Reply to every other student's presentation. Be kind, and be specific: How well did they research and cite their sources? Was the presentation clear and right for their audience? What's one thing they did well, and one thing that would make it better? Comments like ``good job!'' or ``nice work!'' won't get credit.
    \item \textbf{Turn in your final presentation} in the Brightspace assignment, with everything listed under Deliverables. Posting your video to the discussion isn't the same as turning in the assignment. You still have to turn it in.
\end{enumerate}

\section*{Deliverables}
Turn in, in the Brightspace assignment:
\begin{itemize}
    \item Your video, or a link to it
    \item Your slides
    \item A report with your audience, your 5 or more sources with citations, and a CRAAP test checklist or paragraph for each source. Presentations based on unreliable sources won't get credit.
    \item Screenshots of every comment you made on other students' presentations
\end{itemize}

\section*{Resources}
My Final Presentation Topics list
Link: \href{https://www.aholdengouveia.name/InfoSec/labs/finalpresentation-topics.html}{\textbf{\nolinkurl{https://www.aholdengouveia.name/InfoSec/labs/finalpresentation-topics.html}}}

My guide to evaluating sources
Link: \href{https://www.aholdengouveia.name/Howtos/Sourceeval/sourceEvaluation.html}{\textbf{\nolinkurl{https://www.aholdengouveia.name/Howtos/Sourceeval/sourceEvaluation.html}}}

The CRAAP test, with questions to ask about each source
Link: \href{https://web.archive.org/web/20250814000638/https://guides.lib.uchicago.edu/c.php?g=1241077&p=9082343}{\textbf{\nolinkurl{https://web.archive.org/web/20250814000638/https://guides.lib.uchicago.edu/c.php?g=1241077&p=9082343}}}

OBS Studio, free screen recording software for Windows, Mac, and Linux
Link: \href{https://obsproject.com/}{\textbf{\nolinkurl{https://obsproject.com/}}}

Apple: how to record the screen on a Mac
Link: \href{https://support.apple.com/en-us/102618}{\textbf{\nolinkurl{https://support.apple.com/en-us/102618}}}

\end{document}
